% Paragraphs TeX set wider than the column, reaching past its right edge by a fixed measure at every width: a
% list with a negative \rightmargin (as changepage's adjustwidth makes), a negative \rightskip, and a negative
% \leftskip on the left; a centred line in each, and the text at the column's width between them.
\documentclass{article}
\usepackage[T1]{fontenc}
\usepackage{lmodern}
\begin{document}
A justified paragraph at the column's width, long enough to fill a few lines at every width the text is set
at, so that its spaces are stretched or shrunk as TeX chose them and its last line ends in the middle.
\begin{list}{}{\leftmargin 0pt \rightmargin -40pt \listparindent\parindent}
\item[] A justified paragraph in a list whose right margin is negative: its lines reach 40pt past the
column's right edge, as TeX set them, and its spaces are stretched or shrunk at that measure.

\centering Centred in the wider list
\end{list}
{\rightskip -30pt A justified paragraph with a negative right skip, reaching 30pt past the right edge of the
column at every width, its lines broken at that measure and justified to it.\par}
{\leftskip -25pt A justified paragraph with a negative left skip, starting 25pt left of the column and broken
at a measure 25pt wider than the column's, so that it overhangs the left margin on every line.\par}
{\leftskip 0pt plus 1fil \rightskip -40pt plus 1fil \parfillskip 0pt Centred, 20pt right of the column's centre\par}
A last justified paragraph at the column's width, long enough to fill a few lines, after the wide ones.
\end{document}
